\chapter{Symmetry and decoupling}
\label{ch:symmetry}

Identifying all of \eqref{eq:fossen} at once would mean estimating dozens of coefficients
from data in which every axis moves at the same time. The plan in \file{maneuvers.md} is to
identify one axis at a time instead. This chapter explains why that is legitimate and where
it stops being legitimate.

\section{The reflection argument}

POLARIS is (approximately) mirror-symmetric about its $xz$ plane: port and starboard are
mirror images. Consider the reflection $y \mapsto -y$. Under it
\begin{itemize}
  \item the \textbf{longitudinal} variables $u$, $w$, $q$ and forces $X$, $Z$, $M$ are unchanged;
  \item the \textbf{lateral} variables $v$, $p$, $r$ and forces $Y$, $K$, $N$ change sign.
\end{itemize}
A symmetric vehicle cannot tell a manoeuvre from its mirror image, so its equations must be
invariant under the reflection. Consequently:
\begin{keyidea}
A longitudinal force ($X$, $Z$, $M$) can depend on the lateral variables only through
\emph{even} combinations such as $v^2$, $r^2$ or $vr$. A lateral force ($Y$, $K$, $N$)
must be \emph{odd} in the lateral variables, for example $v$, $r$, $ur$ or $uv$.
\end{keyidea}
After linearising about straight-line motion ($v = p = r = 0$), every even combination
vanishes to first order, and the inertia and damping matrices become block-diagonal after
reordering the states:
\begin{equation}
  \vect{M} \;\to\; \begin{bmatrix} \vect{M}_\mathrm{lon} & \vect{0} \\ \vect{0} & \vect{M}_\mathrm{lat} \end{bmatrix},
  \qquad
  \text{lon} = (u, w, q), \quad \text{lat} = (v, p, r).
  \label{eq:blockdiag}
\end{equation}
Surge, heave and pitch form one group, and sway, roll and yaw the other. Exciting surge alone
does not, to first order, produce sway, roll or yaw. Each group can therefore be identified
from its own manoeuvres.

\section{What survives the symmetry}

The decoupling is only \emph{first-order}. Three kinds of coupling remain.
\begin{enumerate}
  \item \textbf{Within a group.} Surge, heave and pitch can still couple to each other, for
        example through off-diagonal added mass $X_{\dot w}$ or $M_{\dot u}$, or through the
        restoring term $(W-B)\sin\theta$.
  \item \textbf{Nonlinear cross-group terms.} The products in \eqref{eq:surge-full} are even in
        the lateral variables, so they are allowed in surge: $-(m-Y_{\dot v})\,vr$. In the
        lateral group the products of a longitudinal and a lateral variable survive, for example the
        sway force $(m-\Xud)\,ur$ and the Munk yaw moment, which is proportional to $(\Xud - Y_{\dot v})\,uv$. This is the
        ``surge--yaw'' coupling of \file{maneuvers.md} section 2.6: turning while moving
        forward.
  \item \textbf{Actuation-level coupling.} Thrusters do not have to respect the hull's symmetry.
        In POLARIS, \code{MOT_1}, the only surge thruster, sits \qty{3}{\centi\meter} above the
        centre of mass. Its line of action therefore produces a pitch moment
        $M = 0.03\,T$, recorded in \file{model.yaml} as the lever \code{0.0300} in the
        \code{pitch} allocation. Because $M$ and $\theta$ are both measured in FLU, where pitch
        is positive nose-down, forward thrust pitches the nose down. In steady state that
        moment is balanced by the restoring torque,
        $\theta_\mathrm{ss} - \theta_0 = 0.03\,T/M_\theta$, so its size depends on the
        still unknown stiffness $M_\theta$. With the dummy BlueROV2 stiffness
        $\qty{1.316}{\newton\meter\per\radian}$ a \qty{20}{\newton} thrust would give
        \qty{0.46}{\radian}. That shows the lever cannot be dismissed until $M_\theta$ has been measured.
\end{enumerate}

\section{Consequences for the plan}

The per-axis plan follows directly from the argument:
\begin{enumerate}
  \item Identify each axis from manoeuvres that excite only that axis. The other axes are
        given to the model as \emph{measured} inputs from the EKF (the \code{dofs:} section of
        \file{model.yaml}).
  \item Validate the per-axis models on holdout runs.
  \item Command two axes together (section 2.6 of \file{maneuvers.md}). If the per-axis models still
        explain the data, coupling is negligible. If a systematic residual remains, fit a
        correction term while leaving the per-axis coefficients untouched.
\end{enumerate}

\begin{revised}[title=Revised 15 October 2026]
Roll is not modelled any more (\cref{sec:no-roll}). The lateral group of the identified model is
therefore $(v, r)$, with $\phi$ and $p$ replayed from the EKF. The symmetry argument is
unaffected: it only states which groups may be identified separately.
\end{revised}

\begin{incode}
The \code{dofs:} block in \file{model.yaml} switches every DOF between \code{modelled} and
\code{measured}. Currently only \code{surge: modelled}; heading, sway, depth and attitude are
replayed from the EKF inside \code{propagate()}. A wrong surge prediction therefore cannot
be blamed on a wrong yaw model. This is error attribution by construction.
\end{incode}

\begin{questions}
\begin{enumerate}
  \item Is a term $X_{vv}\,v^2$ in the surge equation allowed for a port--starboard symmetric
        vehicle? Is $X_v\,v$?
  \item Why does the coupling check include $r^2$ but not $r$ as a natural Fossen term in
        surge? (The code nevertheless tests $r$ as a generic candidate. Why might that still
        be useful?)
  \item Which of the three kinds of surviving coupling would the combined-axes manoeuvre
        (surge and yaw together) excite?
\end{enumerate}
\end{questions}
